\clase{1}{4 de Marzo}{}

\section{Espacios de medida}

\begin{itemize}
    \item Un espacio \(E\) será cualquier conjunto no vacío.
    \item Cualquier colección de subconjuntos de un espacio \(E\) será denominada una \textbf{clase} (de subconjuntos de \(E\)).
\end{itemize}

\begin{definition}[\(\sigma\)-álgebra]
    Una clase \(\Sigma\) de subconjuntos de un espacio \(E\) se dirá una \textbf{\(\sigma\)-álgebra} si cumple que:
	\begin{enumerate}[i)]
        \item \(\emptyset \in \Sigma\) (en particular, \(\Sigma\) es no vacía).
        \item \(A \in \Sigma \Rightarrow A^c \in \Sigma\).
        \item \((A_n)_{n \in \mathbb{N}} \subset \Sigma \Rightarrow \bigcup_{n=1}^{\infty} A_n \in \Sigma\).
    \end{enumerate}
\end{definition}

\begin{definition}[Espacio medible]
    Un \textbf{espacio medible} es cualquier par \((E, \Sigma)\) donde \(E\) es un espacio y \(\Sigma\) es una \(\sigma\)-álgebra (de subconjuntos de \(E\)).
\end{definition}

\begin{definition}[Medida]
    Dado un espacio medible \((E, \Sigma)\), una aplicación \(\mu : \Sigma \to [0, \infty]\) se dirá una \textbf{medida (no negativa)} sobre \((E, \Sigma)\) si cumple que:
	\begin{enumerate}[i)]
        \item \(\mu(\emptyset) = 0\),
        \item si \((A_n)_{n\in\mathbb{N}} \subset \Sigma\) son conjuntos disjuntos entonces \(\mu\left( \bigcup_{n=1}^{\infty} A_n \right) = \sum_{n=1}^{\infty} \mu(A_n)\). 
    \end{enumerate}
    Si \(\mu\) cumple (ii) decimos que \(\mu\) es \textbf{\(\sigma\)-aditiva}.
\end{definition}

\begin{definition}[Espacio de medida]
    Un \textbf{espacio de medida} será cualquier terna \((E, \Sigma, \mu)\) donde \((E, \Sigma)\) es un espacio medible y \(\mu\) es una medida sobre \((E, \Sigma)\).
\end{definition}

\subsection{Ejemplos de medidas}

\begin{eg}[Medida de Lebesgue]~
    \begin{itemize}
        \item $E = \mathbb{R}^d$

        \item $\Sigma = \{\text{conjuntos medibles Lebesgue}\}$

        \item $\mu = \text{medida de Lebesgue en } \mathbb{R}^d$
    \end{itemize}
\end{eg}

\begin{eg}[Medida de contar]~
    \begin{itemize}
        \item $E \text{ un espacio arbitrario}$

        \item $\Sigma = 2^{E} \; (\text{partes de } E)$

		\item $\mu = \text{medida de contar } (\mu(A) = \#A)$
    \end{itemize}
\end{eg}

\begin{eg}[Medidas de Hausdorff]~
    \begin{itemize}
        \item $E = \mathbb{R}^d$

        \item $\Sigma = \{\text{Borelianos de } \mathbb{R}^d\}$

        \item $\text{Fijado } s \in [0,\infty), \quad \mu = \mathcal{H}^s = \text{medida de Hausdorff } s\text{-dimensional}$
    \end{itemize}
\end{eg}

\subsection{Propiedades de una medida}
Si \(\mu\) es una medida sobre \((E, \Sigma)\), entonces:

\begin{property}[Monotonía]
    \(A \subseteq B \implies \mu(A) \leq \mu(B)\).
\end{property}

\begin{property}[Subaditividad]
    \(A \subseteq \bigcup_{n=1}^{\infty} A_n \implies \mu(A) \leq \sum_{n=1}^{\infty} \mu(A_n)\).
\end{property}

\begin{property}[Continuidad inferior]
    \(A_n \nearrow A \implies \mu(A_n) \nearrow \mu(A)\).
\end{property}

\begin{property}[Continuidad superior]
    \(A_n \searrow A\) y \(\mu(A_n) < \infty\) para algún \(n\) (o \(\mu(A_1) < \infty\)) \(\implies \mu(A_n) \searrow \mu(A)\).
    Puede ser falso sin esta condición.
\end{property}

\section{Espacios de probabilidad}

\begin{definition}
    Un espacio de medida \((E, \Sigma, \mu)\) se llamará un \textbf{espacio de probabilidad} si \(\mu(E) = 1\). En este caso, \(\mu\) recibe el nombre de \textbf{medida de probabilidad} (o probabilidad a secas).
\end{definition}

\subsection{Notación para espacios de probabilidad}
\begin{itemize}
    \item El espacio \(E\) se denota \(\Omega\), y lo llamamos \textbf{espacio muestral}.
    \item Los elementos de \(\Omega\) se denotan \(\omega\) (y a veces se llaman \textbf{realizaciones}).
    \item La \(\sigma\)-álgebra \(\Sigma\) se suele denotar por \(\mathcal{F}\).
    \item Los elementos de la \(\sigma\)-álgebra \(\mathcal{F}\) se llaman \textbf{eventos}.
    \item La medida \(\mu\) se denota \(\mathbb{P}\) o \(P\).
\end{itemize}

\subsection{Ejemplos de espacios de probabilidad}

\begin{eg}[Espacios equiprobables]
    Se llaman así porque todas las realizaciones tienen la misma probabilidad.
    \begin{itemize}
        \item \(\Omega\) un conjunto finito.
        \item \(\mathcal{F} = 2^{\Omega}\) (partes de \(\Omega\)).
        \item \(\mathbb{P}(A) = \dfrac{\#A}{\#\Omega}\).
    \end{itemize}
\end{eg}

\begin{eg}[Espacios discretos]
    \begin{itemize}
        \item \(\Omega\) un conjunto finito o numerable.
        \item \(\mathcal{F} = 2^{\Omega}\).
        \item \(\mathbb{P}(A) = \sum_{\omega \in A} p(\omega)\), donde \((p(\omega))_{\omega \in \Omega}\) es algún vector fijo que satisface \(p(\omega) \geq 0 \; \forall \omega \in \Omega\) y \(\sum_{\omega \in \Omega} p(\omega) = 1\).
    \end{itemize}
\end{eg}

\begin{eg}[Un ejemplo no discreto]
    \(\Omega = [0,1]\), \(\mathcal{F} = \{\text{Borelianos de } [0,1]\}\), \(\mathbb{P}(A) = |A|\) (medida de Lebesgue).
\end{eg}

\subsection{Propiedades adicionales de una probabilidad}
Si \(\mathbb{P}\) es una probabilidad sobre \((\Omega, \mathcal{F})\), entonces a las propiedades anteriores se suman:
\begin{property}[Continuidad superior sin hipótesis adicional]
    \(A_n \searrow A \implies \mathbb{P}(A_n) \searrow \mathbb{P}(A)\).
\end{property}
\begin{property}[Complemento]
    \(\mathbb{P}(A^c) = 1 - \mathbb{P}(A)\).
\end{property}

\section{Análisis vs Probabilidad}
Aunque tanto en Análisis (Teoría de la medida) como en Probabilidad se trabaja con espacios de medida, la interpretación en cada área es muy diferente:

\begin{itemize}
    \item \textbf{En Análisis:}
        \begin{itemize}
            \item \(E\) suele ser un espacio topológico (i.e. con alguna noción de geometría dentro), típicamente \(\mathbb{R}^d\).
            \item \(\Sigma\) representa la clase de “conjuntos medibles”, i.e. conjuntos suficientemente buenos geométricamente como para poder “medirles”. 
            \item \textbf{Ejemplo:} la clase de conjuntos medibles Lebesgue en \(\mathbb{R}^d\).
            \item \(\mu(A)\) representa la “medida” del conjunto \(A\). Fijado el espacio \(E\) bajo estudio, típicamente \(\Sigma\) queda fija y \(\mu\) cambia.
        \end{itemize}
    
    \item \textbf{En Probabilidad:}
        \begin{itemize}
            \item \(\Omega\) es el conjunto de posibles resultados de un experimento bajo estudio. No tiene por qué haber geometría subyacente.
            \item \(\mathcal{F}\) representa el conjunto de “información observable” sobre el experimento en cuestión.
            \item \(\mathbb{P}(A)\) representa la probabilidad de que ocurra el evento \(A\). Fijado el experimento bajo estudio, típicamente \(\mathbb{P}\) queda fija y \(\mathcal{F}\) cambia.
        \end{itemize}
\end{itemize}

